\begin{proof}

\begin{enumerate}[label={$\langle1\rangle$\arabic*.},leftmargin=*]
\item Put $N=N_{D/Y}$ and $B=\OO_Y(S)|_D$. Since every ambient line bundle has zero rational first Chern class, $\deg B=\deg K_Y|_D=0$. Adjunction for the smooth curve gives
\[
 \det N=K_D\otimes K_Y^{-1}|_D,\qquad
 \deg N=2g(D)-2.                                      \tag{C.J.2}\hypertarget{tag-00OW}{}\label{cdp:eq:C.J.2}
\]
A local defining section of $S$ belongs to $I_D^2$. Its quadratic normal jet is a well-defined section of $\operatorname{Sym}^2N^*\otimes B$. In the first two cases its generic rank is one. Its saturated radical $R\subset N$ is a line subbundle on the smooth curve, and $Q=N/R$ is a line bundle. The quadratic jet gives a nonzero section of $B\otimes Q^{-2}$, so
$\deg Q\le0$.

\item In the cusp case, the cubic jet on the radical is an intrinsic nonzero section of $B\otimes R^{-3}$. To check its independence, choose coordinates $(x,y,z)$ with $D=(x=y=0)$ and $R$ the $y$ direction. A change of lift of a radical vector changes the cubic by a Hessian term with one radical argument, which is zero. Changing the local retraction to $D$ differentiates the identity that the quadratic jet vanishes on $R$, again giving zero. Thus the local cubics glue across all points, including possible zeros of the quadratic coefficient. Consequently $\deg R\le0$. Equation \eqref{cdp:eq:C.J.2} proves $g(D)\le1$.

\item In the tacnode case the radical cubic vanishes identically. Write the relevant normal terms as
\[
 a x^2+bxy^2+cy^4.
\]
Completing the square has invariant numerator
\[
 \Delta=ac-b^2/4\in
 H^0(D,B^2\otimes Q^{-2}\otimes R^{-4}).               \tag{C.J.3}\hypertarget{tag-00OX}{}\label{cdp:eq:C.J.3}
\]
This numerator is holomorphic even at a zero of $a$. Under $x\mapsto x+uy^2$, the coefficients become $b+2au$ and $c+bu+au^2$, and \eqref{cdp:eq:C.J.3} is unchanged. Linear changes preserving the radical give precisely the displayed line-bundle weights. Changes of the retraction to $D$ add terms containing the radical cubic or a radical Hessian argument, so also leave this quartic numerator unchanged. These checks give its global definition. It is nonzero by the generic tacnode assumption. Hence
\[
 \deg Q\le0,\qquad \deg Q+2\deg R\le0,
 \qquad \deg N\le\tfrac12\deg Q\le0.
\]
This proves the same genus bound.

\item For the ordinary triple case the quadratic jet is identically zero. The binary cubic normal jet is a section of $\operatorname{Sym}^3N^*\otimes B$. Its discriminant is a nonzero section of
$B^4\otimes(\det N^*)^6$. The discriminant has degree four in the cubic coefficients and weight six under a linear change of the two normal coordinates; these weights also follow by writing it as the square of the product of the three pairwise determinants of the linear factors. Its nonnegative divisor degree gives $\deg N\le0$ and the genus bound again.

\item When $g(D)=1$, all the inequalities are equalities. In the cusp case both the quadratic and radical cubic sections have no zeros. In the tacnode case $a$ and $\Delta$ have no zeros. In the triple case the cubic discriminant has no zeros. The holomorphic splitting lemma with parameter $z$ completes the square in the first two cases. The remaining one-variable term is respectively a unit times $y^3$ or $y^4$; taking a local root of the unit gives the corresponding equation in \eqref{cdp:eq:C.J.1}. In the triple case its three distinct tangent lines split after an etale ordering cover; the holomorphic implicit function theorem gives three smooth branches meeting transversely in each normal plane. Straightening two branches to the axes and rescaling the third branch's unit gives the third equation. Thus these are constant transverse types throughout $D$. In particular no other singular curve meets $D$.

\item In the cusp normalization, $x=u^3,y=u^2$, up to the choice of which coordinate is the radical. Its normal line $P=N_{C/\widetilde S}$ satisfies $P^2=R$ and $P^3=Q$, so has degree zero. For the tacnode, the two branches have $x=\pm y^2$ and normal line the pullback of $R$, again of degree zero. Their permutation cover is etale because $\Delta$ is nowhere zero. For the ordinary triple, order the three tangent lines by a finite etale cover. Three pairwise transverse line subbundles in a rank-two bundle make it a sum $P\oplus P$: the third is a graph of an isomorphism between the first two. Since the determinant has degree zero, $\deg P=0$. Each normalization branch has this tangent-root normal line. Etale descent gives the claimed degrees and covers before ordering.

\item Finally the inverse-Jacobian assertion can be checked on these exact transverse equations. For the cusp, $R=\C\{u^2,u^3\}$ and $\mathcal J^{-1}=u^{-1}\C\{u\}$, as in \ref{cdp:st:A.2.4}. For the tacnode, normalization functions are pairs whose constant and linear terms agree. A putative pole of order two in a fractional multiplier has leading coefficients $(a_1,a_2)$. Multiplication by $2x$ forces $a_1=-a_2$ at order zero, while multiplication by $-4y^3$ forces $a_1=a_2$ at order one. Thus both vanish. For the ordinary triple, use branches $(0,u),(u,0),(u,u)$. The partial derivatives have leading coefficient vectors $(-1,0,1)$ and $(0,1,-1)$ at order two. A pole of order two would make both products constant on all three branches, forcing all its leading coefficients to vanish. A pole of larger order is already excluded by regularity of these products on normalization. These computations prove the simple-pole bound. The covariance and leading-coefficient equations are checked exactly in \path{code/C_jets.py}.

\end{enumerate}
\end{proof}
